\section{Cached-force gates and whole-prefix accounting}
\label{app:source-force-screen}

The physical sampler already stores a paired-oracle force at each current state.
We test using it in a cheap gate, with candidate force available only after the
first test passes. Such gates need not be reciprocal. For forward and reverse
passing probabilities $g_f,g_r\in[e^{-B},1]$, the correction is
\begin{equation}
 \alpha_2=\min\{1,e^R g_r/g_f\},\qquad
 \alpha_{\mathrm{total}}=\min\{g_f,e^R g_r\}.
\end{equation}
The reverse gate is evaluated with the newly queried candidate force. The
forward routine receives only the source force. This is ordinary MH on a
screened proposal with rejected mass assigned to self transitions; substituting
$R-s$ without a reciprocal-factor relation would be incorrect. Tests verify
zero-screen RNG/query equivalence, finite-state balance, actual molecular
information order and visible failure after an invalid paid reverse gate.

The five-coefficient control adds source even-force work to the preceding cheap
features. A fixed work control combines this first-order physical energy estimate
with the exact known COM-restraint change and proposal ratio. The neural model
uses a two-layer invariant paired-geometry graph encoder with source force/work,
force norm, displacement, electronic conditions and symmetric leaf/anchor roles.
All 1,671 pairs retain their original split and failures. Forces and projected
scores for 3,233 used states reconstruct exactly from the stored raw/inverted
outputs; no new oracle calls are needed.

We also change the declared proxy to include initial and nonjoint costs and work.
Each recorded physical trajectory contributes equally. Its nonjoint terms remain
fixed while joint expectations are replaced by those of the screen. The 72 FIT
trajectories each cost 128 calls and have mean expected work 0.4947591\,eV, giving
the fixed optimization rate $\eta_0=0.0038653057$\,eV per call. These are still
fixed-source expectations, not a replay of the changed chain. Two seeds per
learned model run 300 steps. Constant thinning is separately calibrated to each
model's FIT joint-query use and evaluated without tuning on selection outcomes.

\begin{table}[h]
\centering\small
\begin{tabular}{lrrr}
\toprule
Gate & Work (eV/prefix) & Calls/prefix & $10^3$ work/call \\
\midrule
No screen & 0.53344 & 128.000 & 4.1675 \\
Fixed physical & 0.50311 & 106.296 & 4.7331 \\
Source work & 0.49293 & 109.003 & 4.5222 \\
Linear, seed 0 & 0.48188 & 114.081 & 4.2241 \\
Linear, seed 1 & 0.49525 & 116.162 & 4.2635 \\
Neural, seed 0 & 0.53344 & 120.684 & 4.4201 \\
Neural, seed 1 & 0.48793 & 114.576 & 4.2585 \\
\bottomrule
\end{tabular}
\caption{Internal whole-prefix expectations on 12 selection parents and two
physical trajectories each. All initial/nonjoint contributions are retained.
The last column is in $10^{-3}$\,eV per expected call. These values are not
achieved speedups or screened-chain endpoint measurements.}
\label{tab:source-force-screen}
\end{table}

All final/control metrics replay. Independent checks cover 6,388 supported-pair
cases, both gates and their balance relation (maximum error
$2.09\times10^{-14}$); six sensitive trained-head finite differences have maximum
error $8.51\times10^{-10}$. Relative to no screening, neural point improvements
are 6.06\% and 2.19\%, versus 13.57\% for fixed physical screening and 8.51\% for
the fixed work control. Both neural-minus-physical 95\% parent intervals are
below zero: $[-5.38,-1.03]\times10^{-4}$ and
$[-8.88,-1.09]\times10^{-4}$\,eV per expected call. Only the first neural seed
establishes improvement over its matched thinning control; neural-minus-linear
intervals span zero for both seeds. These descriptive intervals fix composition
and are not multiplicity-adjusted. This comparison therefore does not establish
a useful neural advantage and does not justify scaling the weights. The earlier
93\% physical-screen result concerns conditional joint-query rates; it is not
inconsistent with the smaller whole-prefix estimate here. Neither estimate
establishes a complete-chain or wall-time gain.
